\documentclass[12pt]{article}
\usepackage{amsmath,amssymb,amsfonts}
\usepackage[margin=1in]{geometry}

\begin{document}

\textbf{Chapter 8, Problem 10.} In Problem 8.61, assume that $V(x) = V_0 =$ const., with $V_0$ small in the sense of Problem 8.5. Calculate $\psi$ and $d\sigma$ using the lowest-order approximation to $\psi$. What is the transmission probability in this approximation? What is the probability of reflection? Note that these two probabilities do not add up to 1. Why not?

(For exact answers, see L.\ I.\ Schiff, \textit{Quantum Mechanics}, McGraw-Hill, second edition, 1955, Section 17.)

\bigskip
\textbf{Solution.}

Consider scattering in one dimension from a constant potential $V(x) = V_0$ for $|x|<a$ and $V(x)=0$ for $|x|>a$. The Lippmann--Schwinger equation is:
\[
\psi(x) = e^{ikx} + \int_{-a}^{a} G_0(x,x')V_0\,\psi(x')\,dx'.
\]

\textbf{First Born approximation:} Replace $\psi(x')$ by $e^{ikx'}$ in the integral:
\[
\psi^{(1)}(x) = e^{ikx} + V_0\int_{-a}^{a} G_0(x,x')e^{ikx'}\,dx'.
\]

Using $G_0(x,x') = \frac{im}{\hbar^2 k}e^{ik|x-x'|}$:

\textbf{Transmitted wave} ($x > a$): Here $|x-x'| = x-x'$ for $x'>-a$, so:
\[
\psi^{(1)}(x) = e^{ikx} + \frac{imV_0}{\hbar^2 k}\int_{-a}^{a} e^{ik(x-x')}e^{ikx'}\,dx' = e^{ikx}\left[1+\frac{imV_0}{\hbar^2 k}\cdot 2a\right].
\]
Thus the transmission amplitude is:
\[
T \approx 1 + \frac{2imaV_0}{\hbar^2 k}.
\]
The transmission probability:
\[
|T|^2 \approx 1 + \frac{4m^2a^2V_0^2}{\hbar^4 k^2}.
\]

\textbf{Reflected wave} ($x < -a$): Here $|x-x'| = x'-x$, so:
\[
\psi_{\text{refl}}(x) = \frac{imV_0}{\hbar^2 k}\int_{-a}^{a}e^{-ikx}e^{2ikx'}\,dx' = \frac{imV_0}{\hbar^2 k}e^{-ikx}\cdot\frac{e^{2ika}-e^{-2ika}}{2ik} = \frac{mV_0}{\hbar^2 k}\frac{\sin(2ka)}{k}e^{-ikx}.
\]
The reflection amplitude:
\[
R = \frac{mV_0\sin(2ka)}{\hbar^2 k^2}.
\]
The reflection probability:
\[
|R|^2 = \frac{m^2V_0^2\sin^2(2ka)}{\hbar^4 k^4}.
\]

\textbf{Why don't they add to 1?}
\[
|T|^2 + |R|^2 = 1 + \frac{4m^2a^2V_0^2}{\hbar^4 k^2} + \frac{m^2V_0^2\sin^2(2ka)}{\hbar^4 k^4} \neq 1.
\]

The Born approximation violates unitarity because it is only the first-order term in a perturbation expansion. Conservation of probability ($|T|^2+|R|^2=1$) is an exact relation that requires summing the full Born series. At each finite order of perturbation theory, unitarity is only satisfied to the corresponding order in the small parameter $mV_0 a/(\hbar^2 k)$. To first order in $V_0$, we have $|T|^2\approx 1$ and $|R|^2\sim V_0^2$ (second order), so they are consistent to order $V_0$ but not exactly. $\blacksquare$

\end{document}
